\cotefile{64}
\scanpage{53}\mspage{5}
\[
 a)\longrightarrow b)\qquad \Gamma=\underline{\mathrm{Norm}}_G\,\underline V
\]
$b)\longrightarrow a)$ \quad $\underline V$ est l'unique sous-schéma en groupes \inserted{distingué} fini étale de $\Gamma$, \deleted{\unsure{g}de} degré $4$, \deleted{distingué}.
\margin{\uline{NB} $\Gamma/V\simeq\underline{\mathrm{Aut}}_{S\mathrm{gr}}(V)\simeq\underline{\mathrm{Aut}}_S(V^*)$}
\[
 c)\longrightarrow b)\qquad \Gamma=\underline{\mathrm{Aut}}_S(X,\Delta)\xleftarrow{\sim}\underline{\mathrm{Aut}}^+_{S\text{-oct}}(\Delta)
\]
$a)\longrightarrow c)$ \quad \deleted{$\underline V^*$ est} Toute section $\sigma$ de $V^*=V\smallsetminus(\text{section unité})$ définit une action \inserted{fidèle} de \deleted{\unsure{ce|st}} $(\mathbb Z/2\mathbb Z)_S$ sur $X$, donc un ``axe'' $\Delta_\sigma$ de $X$, et $\Delta$ est \deleted{la réunion} (loc. ét.) le \deleted{\unsure{sché}} diviseur somme de ses axes \ldots

\uline{Corollaire} 1 \quad Pour un diviseur étale de degré 6 $\Delta$ dans $X$, il existe au plus une structure icosaédrale régulière~; \deleted{\unsure{dont} le \inserted{groupe} \illegible{1} des autom.} \inserted{S'il en existe une, le} \inserted{schéma en} groupes de ses automorphismes \inserted{directs} est isom. à $\underline{\mathrm{Aut}}_S^{\deleted{\unsure{+|4}}}(X,\Delta)$.

\uline{Corollaire 2} \quad La catégorie des couples $(X,\Delta)$ \inserted{\uline{sur $S$}}, (ou $(X,\underline V)$, $(X,\underline\Gamma)$ \ldots) \deleted{\unsure{orientés}} \inserted{est} équivalente à celle des octaèdres \deleted{\unsure{orientés}} combinatoires \inserted{$\Delta$} \inserted{sur} $S$, ou encore à celle des revêtements finis étales \deleted{$\unsure{X|V|\underline V}$} \inserted{$\widetilde F$} de degré 4 de $S$, en associant à \unsure{tt} $(X,\Delta)$ le revêtement des quatre ``faces modulo antipodisme'' de la structure octaédrale $\Delta$, ou ce qui revient au $=$, celui des \inserted{3-}sous-groupes de Sylow de \deleted{$\unsure{G|\underline{\mathfrak S}}_{\Delta}$} $\underline{\mathrm{Aut}}^+_{S\text{-oct}}(\Delta)\simeq\underline{\mathrm{Aut}}_S(X,\Delta)$.
\inserted{munis d'un iso $\underline{\unsure{\mathrm{Or}|\mathrm{or}}}\,\Delta/S\simeq\mu^*_{4S}$ (on dit qu'ils sont $\mu^*_4$-orientés)}
La structure-type $*$ correspond à \deleted{\unsure{inscrire à|brouillon|revêt.}} $X\subset\mathbb P^2_{\mathbb Z[\frac12]}$, $X=\mathbb V(x^2+y^2+z^2=0)$, avec l'op. de \deleted{$\unsure{V}$} $\{\pm\}^3/\mathrm{diag}\simeq\mathbb V_{\{0,1,\infty\}}(\mathbb F_2)$ donnée par les chgts de signe en $x,y,z$ $*$.
% margin note written sideways
\margin{\uline{NB} Ici $X$ est totalement anisotrope sur $\mathbb R$ i.e. n'a pas de pts réels sur $\mathbb R$ -- c'est le revêtement galoisien \unsure{universel} \unsure{absolument} de groupe $\mathbb V_{\{0,1,\infty\}}(\mathbb F_2)$ de $\mathbb P^1_{\deleted{\unsure{\mathbb Z}}}$, ramifié \unsure{seulem\textsuperscript{t}|seulement} en $0,1,\infty$ \unsure{(en} \ldots}
% margin note written sideways
\margin{$*$ La \inserted{bi}face distinguée est formée des deux faces $\{(0,1,i),(i,0,1),(1,i,0)\}$ et $\{\{0,1,-i\},\{-i,0,1\},\{1,-i,0\}\}$}
% archive sheet number: 52

\scanpage{54}
2. \uline{Courbes relatives de type} $(0,4)$.

Soit $X$ un fibré en droites projectives sur $S$, et
\[
 \underline I\subset X \tag*{(9)}
\]
un diviseur relatif, étale de degré 4. Considérons
\[
 \underline{\mathfrak S}_{\underline I}=\underline{\mathrm{Aut}}_S(I) \tag*{(10)}
\]
qui est un schéma en groupes fini étale sur $S$, loc isom. à $(\mathfrak S_4)_S$. Il lui correspond un unique sous-groupe \inserted{distingué} $\underline V_{\underline I}$, loc. isom. à $(\unsure{\mathbb D_{2v}|\mathbb D_{22}|\mathbb D_2})_S$, \deleted{de telle \unsure{faç|façon}} Alors $\underline I$ est un torseur sous $\underline V_I$.

\uline{Proposition} \quad L'opération de $\underline V_{\underline I}$ sur $\underline I$ se prolonge (de façon évidemment unique) en une opération de $\underline V_{\underline I}$ sur $X$.

\uline{NB} La donnée de $\underline I$, fini étale de degré 4 sur $S$, équivaut à celle de $\underline V=\underline V_{\underline I}$ sur $S$ (ou encore de $\underline V_{\underline I}^*=\underline J_{\underline I}$, fini étale de degré 3 sur $S$), plus la donnée d'un torseur $\underline I$ sous $\underline V^*$. Plonger alors $\underline I$ dans un fibré en droites projectives $X$, revient à faire opérer fidèlement $\underline V$ sur $X$ (chose étudiée \unsure{en|au} § 1), et de plus, de se
% archive sheet number: 53

\scanpage{55}\mspage{6}
donner une section de $(X\smallsetminus\Delta)/\underline V$ \deleted{\unsure{de}}, en prenant \inserted{pour $I$} l'image inverse de cette section dans $X\smallsetminus\Delta$.

Mais la donnée d'une structure octaédrale équivaut à celle d'un \deleted{$\unsure{\underline I|\underline J}$} rev. étale de degré quatre de $S$ (à ne pas confondre avec le $\underline I$ de tantôt, appelons le $\underline K$), ou encore, à celle de

a) un schéma en groupes $V$ sur $S$, loc. isom. à $(\mathbb D_2)_S$ -- ou encore un revêtement fini étale $\underline J$ ($=\underline V^*$) de degré 3 de $S$, ou \deleted{de} enfin, à celui d'un \deleted{\unsure{gr}} fibré en droites projectives $\Sigma=\Sigma_{\underline J}$ sur $S$, muni d'un diviseur \inserted{relatif} étale $D_J$ de degré 3 (\deleted{\unsure{si}} isom. \deleted{\unsure{lui}} \inserted{à} $\underline J$)

b) d'un torseur $\underline K$ sous $\underline V$, ou ce qui revient au $=$, d'un revêtement \deleted{\unsure{galoisien}} \inserted{principal} \inserted{$X^*$} de $\Sigma_{\underline J}^*=\Sigma_{\underline J}\smallsetminus\deleted{\unsure{J}}D_{\underline J}$ de groupe $\underline V$, \deleted{\unsure{géométriquement connexe sur chaque fibre}} correspondant à \deleted{$\unsure{\xi}$} un élément de \deleted{$\unsure{H^0(S,R^1 f_* V)|H^0(S,R^1 f_*(V)}$} $H^0(S,\mathcal H^1(\Sigma_{\underline J}^*/S,\underline V))\simeq\mathrm{Hom}(\underline V_{\underline J},\underline V)$ \deleted{$\unsure{=\underline V_{\underline J}|/\mathrm{iso}}$} qui est l'identité $\underline V_{\underline J}\longrightarrow \underline V=\underline V_{\underline J}$.
% slanted
\margin{\uline{NB} Le syst. local des $(\pi_1)_{\mathrm{ab}/2}$ de $\Sigma_{\underline J}^*$ sur $S$ est \deleted{$\unsure{V}$} can. isom. à $V_{\underline J}$ \ldots}
-- La correspondance entre $\underline K$ et $X$ étant obtenue \deleted{par} la formule
\[
 X^*=\deleted{\unsure{X_{J,0}\wedge_{\Gamma}|\Sigma_{\underline J,0}\wedge_{\Gamma}}}\ K\mathbin{\wedge}_{V}\unsure{\overline\Sigma|\Sigma}_{\underline J,2}^*, \tag*{(11)}
\]
où $\overline\Sigma_{\underline J,2}^*$ est le revêtement principal \unsure{dié|dû|dé|dit}
% archive sheet number: 54

\scanpage{56}
% archive sheet number: 55
$\Sigma^*_{\underline J}$ défini \deleted{par} à partir de \deleted{la fibré} la cônique \deleted{\unsure{déf}} d'équation homogène
\[
 \sum_{i\in J}\unsure{y|x}_i^2=0 \tag*{(12)}
\]
% slanted
\margin{description valable loc\textsuperscript{t}. (ét)~; sinon, il faut prendre, dans $\check{\mathbb P}(p_*(\mathcal O_{\underline J}))$, la conique d'équation \deleted{\unsure{de}} $Q=0$, où $Q(u\deleted{\unsure{x|\otimes}})=\mathrm{Tr}\,u^2$ (NB $p:\underline J\longrightarrow S$ le morphisme structural)}
sur laquelle le groupe $V_{\underline J}$ \inserted{$\simeq\mathbb F_2^{\underline J}/\mathbb F_2$} opère \deleted{par un} \inserted{par} changement de signes sur les $x_i$, avec comme quotient $\Sigma_{\underline J}$, défini par l'équation \deleted{homogène}
\[
 \sum x_i=0 . \tag*{(13)}
\]
% slanted
\margin{description valable loc. ét. -- la description intrinsèque est \deleted{\unsure{par}} l'hyperplan (13) de $\check{\mathbb P}(p_*(\mathcal O_{\underline J}))$, d'équation $L\deleted{\unsure{(u)}}=0$, où $L(u)=\mathrm{Tr}\,u$}

\deleted{Donc} La donnée de $(X,\underline I)$ équivaut à celle de \deleted{$\unsure{\Sigma}$} $\underline J$ i.e de $\Sigma_{\underline J}$, du torseur $\underline K$ sous $V_{\underline J}$ i.e. du revêtement \deleted{\unsure{en}} $X$ de $\Sigma_{\underline J}$ \deleted{ramifié} de groupe $V_{\underline J}$ (admissible au sens qu'il définit l'\uline{identité} $V_{\underline J}\longrightarrow V_{\underline J}$), \uline{et} d'une section $\unsure{t}$ de $\Sigma_{\underline J}^*$.
% slanted
\margin{Relation entre $\underline I,\underline K,x$~: $(14)\;\big|\;\underline I\simeq\underline K\mathbin{\wedge}_{\underline V}(\overline\Sigma^*_{\underline J,2}|_{t})$. Ainsi, pour $\underline J$ et $x\in\Gamma(\Sigma^*_{\underline J}/S)$ fixés, $\underline I$ et $\underline K$ (ou encore, $\underline I$ et $X$) se déterminent mutuellement \ldots{} On se gardera de confondre les trois rev\textsuperscript{ts} étales de degré 4 de $S$, $\underline I$, $\underline J\amalg S$, et $\underline K$ \ldots}

Quand $\underline I$ est \inserted{fixé et} de la forme $I_S$, \inserted{[l'hypothèse que $\underline I$ soit constant est en fait superflue]} avec $I$ ens. à quatre éléments fixé (à fins d'épinglage), on a $\underline J=J_S$, et $\Sigma_{\underline J}^*=\overline\Sigma_J$ est aussi fixé. La donnée de $(X,I_S\hookrightarrow X)$ équivaut alors à celle de la section $x$ de $\Sigma_J^*$, servant de ``pt référence'', le revêtement $X$ de $\Sigma_J^*$ de groupe $(V_J)_S=(V)_{(J_S)}=\underline V$ étant

\scanpage{57}\mspage{7}
déduisant par la \deleted{\unsure{f}} structure supplémentaire
\[
 X|t\xrightarrow{\sim}I_S\qquad(V_S\text{-isomorphisme}) \tag*{(15)}
\]
qui, en termes du $V_S$-torseur $\underline K$ qui est censé définir $X$, s'écrit
\[
 \deleted{\underline K\simeq(I_S)\mathbin{\wedge}_{\unsure{(V_J)_S|V_S}}(\overline\Sigma^*_{J,2}|t)} \tag*{\deleted{(16)}}
\]
\[
 \underline K\simeq(I_S)\mathbin{\wedge}_{\underline V}(\Sigma^*_{J,2}|t)\qquad \underline V=(V_J)_S. \tag*{(16)}
\]
% archive sheet number: 56

\scanpage{58}
3: \quad Relations avec courbes elliptiques\inserted{~: le morphisme $f_2^t:\Sigma_{\underline J}\longrightarrow\Sigma_{\underline J}$}.

On suppose tjrs les car. rés. de $S\ne 2$. Soient $E_0$ une courbe elliptique sur $S$, $E$ un torseur sous $E_0$ (de sorte que $E_0\simeq\underline{\mathrm{Pic}}^0_{E/S}$).

On suppose donné dans $E$ un \inserted{$S$-}\unsure{autom|automorphisme}
\[
 \sigma:E\longrightarrow E \tag*{(17)}
\]
tel que
\[
 \sigma_{E_0}^{\unsure{2}}=-\mathrm{id}_{E_0} \tag*{(18)}
\]
ce qui implique que loc. (ét), identifiant $E$ à $E_0$ via une section origine de $E$, on aura
\[
 \sigma x=a-x\qquad a\ \text{section convenable}
\]
ce qui est d'ailleurs intrinsèque si on interprète $a$ comme une section de $E\mathbin{\wedge}_{E_0}E\simeq\underline{\mathrm{Pic}}^2_{E/S}$.

On aura donc
\[
 \sigma^2=\mathrm{id}_E
\]
Soit \deleted{$\unsure{\underline I|4}$}
\[
 \left\{\begin{aligned}
 &\underline I=E^\sigma\quad\text{sous-schéma des pts fixes de }\sigma\\
 &\underline I\subset E\\
 &\ \Vert\\[-3pt]
 &\{x\in E\mid 2x=a\}
 \end{aligned}\right. \tag*{(19)}
\]
on voit que $I$ est un torseur sous le schéma en groupes
\[
 \underline V={}_2E_0=\mathrm{Ker}\,(2\,\mathrm{id}_{E_0}:E_0\unsure{\text{—}|\longrightarrow} E_0), \tag*{(20)}
\]
qui est justement un groupe \inserted{sur $S$} loc. isom. à $\mathbb D_{2S}$.

La donnée de $(E,\sigma)$, où $E$ est un schéma
% archive sheet number: 57

\scanpage{59}\mspage{8}
propre et lisse sur $S$ à fibres géom. connexes de dim 1, genre 1, et $\sigma$ un autom induisant $-\mathrm{id}$ sur $\underline{\mathrm{Pic}}^0_{E/S}$ \uline{équivaut} à celle \inserted{de $(E,a)$, où $a$ section de $\underline{\mathrm{Pic}}^2_{E/S}$ ($\simeq E\mathbin{\wedge}_{E_0}E$), ou encore à celle} \deleted{de} $(E_0,\underline I)$, où $E_0$ est une courbe elliptique relative sur $S$, et où $I$ est un torseur sous \deleted{$\,{}_2E_0\overset{\mathrm{déf}}{=}$} $\underline V\overset{\mathrm{déf}}{=}{}_2(E_0)$. On aura
\[
 \left\{\begin{aligned}
 &E=\underline I\mathbin{\wedge}_{\underline V}E_0\\
 &\sigma\ \text{déduit de la symétrie }x\mapsto-x\ \text{de }E_0\\
 &\hspace{2em}\text{via la formule précédente}
 \end{aligned}\right. \tag*{(21)}
\]
On aura une opération de $\underline V\deleted{\subset}={}_2E_0\subset E_0$ sur $E$, et un isom can.
\[
 \begin{aligned}
 &E/\underline V\qquad\simeq\qquad E_0\\
 &\qquad\wr\!|\\
 &E\mathbin{\wedge}_{E_0}E\simeq\underline{\mathrm{Pic}}^2_{E/S}
 \end{aligned} \tag*{(22)}
\]
% slanted
\margin{\uline{NB} L'opération de $\underline V$ sur $E$ commute à l'action de $\sigma$}
isom qui commute aux symétries $\sigma$ sur $E$, $\sigma_0:x\mapsto-x$ sur $E_0$.

Considérons les schémas quotients
\[
 X=E/\{1,\sigma\},\qquad\Sigma=E_0/\{1,\sigma_0\} \tag*{(23)}
\]
ce sont des fibrés en \deleted{\unsure{droites}} \inserted{droites} projectives, \deleted{et on a les} $V$ opère \inserted{(monomorphiquement)} sur $X$ \deleted{et sur $\Sigma$} via son \deleted{\unsure{(24)}} action sur $E$, et on aura
\[
 X/V\simeq\Sigma, \tag*{(24)}
\]
d'ailleurs le diviseur de ramification pour
\[
 E\longrightarrow X\qquad(\text{rev. quadratique}) \tag*{(25)}
\]
% archive sheet number: 58

\scanpage{60}
est $I$, et on a un \deleted{\unsure{mono}} \uline{morphisme}
\[
 I\hookrightarrow X \tag*{(26)}
\]
\deleted{stable \unsure{avec|sous} l'acti} compatible avec l'action de $\underline V$, d'où on déduit facilement que c'est un \uline{monomorphisme} (car on vérifie que $I$ s'envoie dans l'ouvert \deleted{\unsure{o}} $\underset{(27)}{X^*}$ de $X$ où l'action de $\underline V$ est libre, cf plus bas).

Ainsi on est dans la situation du § 2, qui implique que $\Sigma\simeq\Sigma_{\underline J}$, où $\underline J=\underline V^*$, que $I$ est l'image inverse d'une section bien déterminée $t$ de $\Sigma^*_{\underline J}=\Sigma_{\underline J}\smallsetminus D_{\underline J}$, l'image inverse $\Delta$ de $\underline J$ dans $X$ étant \deleted{\unsure{une}} la configuration \unsure{icosaédrale|icosaèdrale} dans $X$, canoniquement associée à l'action donnée de $\underline V$ sur \deleted{le} fibré en droites projectives $\underline V$ sur $X$.

En termes de l'action de $\underline{\deleted{\Delta}V}$ sur $X$, $\underline\Delta$ se récupère comme le diviseur de ramification pour cette action -- ensemblistement, il correspond donc aux pts (à val. dans des corps alg. clos) dont le stabilisateur ds $V$ est $\ne 0$, i.e. il correspond à des pts $x\in E$ tels qu'il existe $\unsure{\zeta|\xi}\in V={}_{\unsure{2|8}}E_0$, $\zeta\ne0$, avec $x+\zeta\in\{x,\sigma x=a-x\}$ i.e $x+\zeta=a-x$ i.e. $2x+\zeta=a$ d'où $4x=2a$ \deleted{i.e. $x\in \unsure{I\mathbin{\wedge}{}_4E_0}$,} \inserted{mais $2x\ne a$}
% archive sheet number: 59

\scanpage{61}\mspage{9}
$x$ pt de $I\mathbin{\wedge}_{{}_2E_0}({}_4E_0)\smallsetminus I$

On trouve ainsi
\[
\Delta=\underbrace{\bigl(\underbrace{\underbrace{I\mathbin{\wedge}_{\underline M_0}\underline M}_{\text{degré }16}
\smallsetminus\underbrace{I}_{\text{degré }4}}_{\text{degré }12}\bigr)/\{1,\sigma\}}_{\text{degré }2}
\tag*{(27)}
\]
où on a posé maintenant
\[
\left\{\begin{aligned}
\underline M&\simeq {}_4E_0\quad\text{schéma en \unsure{groupe|groupes} loc. isom.}\\
&\hspace{5em}\text{à }(\mathbb Z/4\mathbb Z)^2_S\\
\underline M_0&\simeq\underline M\otimes_{\mathbb Z/4\mathbb Z}\mathbb Z/2
\simeq2\underline M\simeq{}_2E_0=\underline V
\end{aligned}\right.
\tag*{(28)}
\]
Ainsi \inserted{(dans le cas $E=E_0$, disons)} les pts d'ordre 4 de $E_0$ se groupent par \inserted{6} paires d'éléments symétriques, qui forment la configuration octaédrale $\Delta$ dans $X$.

On trouve ainsi le diagramme (presque cartésien)
\diagram{\begin{tikzcd}[ampersand replacement=\&,column sep={2.5em,between origins},row sep=1.5em]
{}\&\underline I\arrow[r,phantom,"\subset" description]\& E\arrow[rrrr,"\text{étale}"]\arrow[d,"{\substack{\text{ramifié au}\\\deleted{\unsure{I}}\text{dessus des pts}\\\text{de }\underline I}}"] \&{}\&{}\&{}\& E_0\;(\simeq E/\underline M_0)\arrow[d,"{\substack{\text{ramifié au dessus}\\\text{des pts de }D_J\cup t(S)}}"]\\
{(29)}\&\underline I\arrow[r,hook]\& X\arrow[rrrr,"{\text{ramifié au}}","{\substack{\text{dessus des}\\\text{pts de }D_J}}"']\&{}\&{}\&{}\&\Sigma_J\\[-1.4em]
{}\&{}\&\cup\&{}\&{}\&{}\&\cup\\[-1.4em]
{}\&{}\&\Delta\arrow[rrr,dash,end anchor=center]\&{}\&{}\&{}\arrow[r,dashed,start anchor=center]\&D_J
\end{tikzcd}}
qui est cartésien au dessus de $\unsure{\Sigma|\overline\Sigma}_J^*=\overline\Sigma_J\smallsetminus D_J$
% archive sheet number: 60

\scanpage{62}
(au dessus \unsure{au quel|duquel} $X^*$ est étale \deleted{galoi} principal de groupe $\underline V$), tout comme $E$ sur $E_0$) mais non au voisinage des pts de $D_J\subset\Sigma_J$ (vu que $E_0/\Sigma_J$ y est ramifié, mais non $E/X$ en les points de $\Delta=X\mid D_J$, or si c'était cartésien il en résulterait que $E_0/\Sigma_J$ est étale \unsure{sur|en|aux} les pts de $D_J$ aussi, par descente plate\ldots). Quand $=$ il doit être facile de vérifier que le diagramme total se reconstitue \unsure{à|:} partir de $E_0$ ($\longmapsto$ \inserted{$\underline J$ et} $\Sigma=\Sigma_{\underline J}$), et $E_0\longrightarrow\Sigma$, \inserted{et section $t$ de $\Sigma^*$}) et du revêtement \deleted{principal} $X$ de $\Sigma$ avec groupe \inserted{\unsure{de type|d'op.|d'op|étales}} $\underline V$, : la donnée de $(E,\sigma)$ i.e de $(E_0,I)$ \deleted{doit être} \inserted{est} équivalente aussi à $E_0$ et du revêtement $X$ (d'où en effet \deleted{$\underline I$}$\underline K$ et $\underline I\simeq\underline K\wedge(X\mid t)^{\unsure{\prime|,}}$\ldots)

Cela montre que finalement $(E,\sigma)$ i.e. le $I$ est un peu de la poudre aux yeux, on va prendre par la suite $E=E_0$, $\sigma=\sigma_0$, donc $I=\underline V={}_2E_0=M_0$, torseur trivial. On a donc \deleted{$\unsure{\Sigma|J}$} $X=\Sigma_J$ et on trouve donc
% archive sheet number: 61

\scanpage{63}\mspage{10}
une application canonique
\[
\Sigma_{\underline J}\xrightarrow{\ f^{t}_{\unsure{2|2_0}}\ }\Sigma_{\underline J}
\tag*{(30)}
\]
\margin{c'est l'application déduite par passage aux quotients de $2\operatorname{id}_{E_0}:E_0\longrightarrow E_0$}
% slanted
\margin{il faudrait revenir sur une description intrinsèque de $f_{\unsure{i|2}}^t$, sans faire intervenir \unsure{de courbe elliptique|des courbes elliptiques|de courbes elliptiques}}
qui doit pouvoir se définir en termes de $\underline J$, et de la seule section $t$ de $\Sigma_{\underline J}$ (ce qui en effet détermine $E_0$, à isom \inserted{(non unique)} près loc. pour la top. étale\ldots) On a ici
\[
\deleted{I=J\amalg}\qquad \underline I_t=\underline J\amalg t(S)
\tag*{(31)}
\]
\margin{\uline{NB} comme $S$-schéma, $\underline I_t$ ne dépend pas de $t$, mod iso. can.}
et \deleted{$\unsure{i|I|E_0}$} $\underline V$ peut se décrire comme $V_{\underline I}\simeq V_{\underline J}$, opérant sur $\underline I$ de la façon canonique, cette opération se prolongeant en une opération sur $\Sigma_{\underline J}$ (dépendant de la section $t$ de $\Sigma_{\underline J}^*$). Le \deleted{applic} morphisme $f_t$ est défini sans ambiguïté par les conditions suivantes (où on suppose pour simplifier que $\underline J=J_S$, $J$ ens à trois éléments) :
\[
\left\{\begin{aligned}
&\forall i\in J,\ \text{si }\sigma_i^t\in\operatorname{Aut}_S(\Sigma_J)
  \text{ est l'autom. qui échange}\\
&\qquad t\text{ et }i,\text{ et }j\text{ et }k\quad(\text{où }\{i,j,k\}=J),\text{ et}\\
&\deleted{\text{alors }\unsure{\text{si}|\text{soit}|f}\,(\Sigma_J)^{\sigma_i^t}}
  \qquad\Delta_i^t=(\Sigma_J)^{\sigma_i^t},\ \text{on a}\\
&\hspace{8em}f(\Delta_i^t)\subset\text{section }\zeta_i\text{ de }D_{\underline J}\simeq J_S
\end{aligned}\right.
\tag*{(32)}
\]
et cela doit impliquer automatiquement
\[
f_t(I_S=J_S\cup t(S))\subset\deleted{\text{section}}\;t(S)
\tag*{(33)}
\]
% archive sheet number: 62

\scanpage{64}
Faisons les calculs pour $J=\{0,1,\infty\}\subset\mathbb P^1_S$, on trouve facilement :
\[
\left\{\begin{aligned}
\sigma_0^t(z)&=\frac{z-t}{z-1} &&(\text{échange }0\text{ et }t,\ 1\text{ et }\infty)\\
\sigma_1^t(z)&=t/z &&(\text{échange }1\text{ et }t,\ 0\text{ et }\infty)\\
\sigma_\infty^t(z)&=t\frac{z-1}{z-t} &&(\text{échange }\infty\text{ et }t,\ 0\text{ et }1)
\end{aligned}\right.
\tag*{(34)}
\]
\margin{NB $\sigma_\infty^t=\sigma_0^t\sigma_1^t$}

On trouve donc
\[
\left\{\begin{aligned}
\Delta_0^t &: \text{ solutions de }z^2-2z+t=0,\quad &\Delta_0^t&=\{1\pm\sqrt{1-t}\}\\
\Delta_1^t &: \rule{3em}{.35pt}\quad z^2-t=0 &\Delta_1^t&=\{\pm\sqrt t\}\\
\Delta_\infty^t &: \rule{3em}{.35pt}\quad z^2-2tz+t=0 &\Delta_\infty^t&=\{t\pm\sqrt{t(1-t)}\}
\end{aligned}\right.
\tag*{(35)}
\]
et on doit avoir
\[
\begin{aligned}
f_2^t:\quad &\Delta_0^t\longrightarrow 0\\
&\Delta_1^t\unsure{\text{---}|\longrightarrow}1\\
&\Delta_\infty^t\longrightarrow\infty
\end{aligned}
\qquad\text{\uline{avec multiplicité 2}}
\tag*{(36)}
\]
ce qui nous donne (puisque \deleted{$\unsure{f_2|f}$}les zéros et pôles de $f_2^t$ de degré quatre sont connus)
\[
f_2^t(z)=c_t\left(\frac{z^2-2z+t}{z^2-2tz+t}\right)^2
\qquad c_t\ \text{section de }\underline O_S^*,
\]
dont la valeur s'obtient en écrivant que $\Delta_1^t\longrightarrow1$. Mais si $\unsure{\xi|\zeta}$ est une section de $\Delta_1^t$ i.e. $\xi^2=t$, on a
\[
\begin{aligned}
\xi^2\deleted{+t}-2\xi+t&=2(t-\xi)\\
\xi^2-2t\xi+t&=2t(1-\xi)
\end{aligned}
\]
\[
f_2^t(\xi)=c_t\left(\frac{2(t-\xi)}{2t(1-\xi)}\right)^2
=c_t/t^2\,\deleted{\frac{\xi^2-2t\xi+t^2}{\xi^2-2\xi+1}=}\left(\frac{t-\xi}{1-\xi}\right)^2
\]
Or
\[
\begin{aligned}
\deleted{=\frac{c_t}{t^2}\frac{t^2+t-2t\xi}{1+t-2\xi}}-\quad
\frac{t-\xi}{1-\xi}&=\frac{(t-\xi)(1+\xi)}{(1-\xi)(1+\xi)}
=\frac{-\xi^2+(t-1)\xi+t}{1-\xi^2}\\
&=\frac{(t-1)\xi}{1-t}=-\xi
\end{aligned}
\]
donc
\[
 f_2^t(\xi)=\frac{c_t}{t^2}\xi^2=\frac{c_t t}{t^2}=\frac{c_t}{t}
 \qquad\text{d'où}\qquad c_t=t
\]
% archive sheet number: 63

\scanpage{65}\mspage{11}
et par suite
\[
 f_2^t(z)=t\left(\frac{z^2-2z+t}{z^2-2tz+t}\right)^2
 \tag*{(37)}
\]
Il reste à vérifier que les valeurs de cette fonction pour $z=0,1,\infty$ ou $t$, sont égales à $t$, i.e que celle de $\frac{z^2-2z+t}{z^2-2tz+t}$ est égale à $\pm1$. Pour $0,\infty,1$ on trouve $1$, pour $z=t$ on trouve $\frac{t^2-t}{t-t^2}=-1$, O.K.

\uline{NB} En prenant le $n$.\unsure{ème|ième} itéré de $f_2^t$, on trouve le morphisme $f_{2^n}^t$, qui correspond par passage au quotient à $2^n\operatorname{id}_{E_0}$. L'image inverse de la section $t$ par $f_{2^n}^t=(f_2^t)^{\circ n}$ est l'image par $p:E_0\longrightarrow\Sigma_J\simeq E_0/\{1,\sigma_{E_0}\}$ du \deleted{\unsure{l'ens|l'ense|l'ensemble}} \inserted{sous-}schéma \inserted{en groupes} $\;{}_{2^n}E_0$ des points annulés par $2^n$.
:
% archive sheet number: 64
